\documentclass[a4paper,12pt]{article}
\usepackage[utf8]{inputenc}
\usepackage[T1]{fontenc}
\usepackage[ngerman]{babel}
\usepackage{amsmath, amssymb}
\usepackage{geometry}
\usepackage{tikz}
\usetikzlibrary{shapes.geometric, arrows.meta, positioning}
\usepackage{parskip}
\usepackage{titlesec}
\usepackage{enumitem}

\geometry{left=3cm, right=3cm, top=2.5cm, bottom=2.5cm}
\titleformat{\section}{\large\bfseries}{}{0em}{}[\titlerule]
\titleformat{\subsection}{\normalsize\bfseries}{}{0em}{}

\tikzset{
  box/.style={rectangle, draw, rounded corners=3pt, text width=4.8cm,
              align=center, minimum height=1.0cm, font=\small, inner sep=4pt},
  arrow/.style={-{Latex[length=2mm]}, thick},
}

\begin{document}

\begin{center}
  {\LARGE\bfseries Studierhinweise zu Lektion 6}\\[0.5em]
  {\large Lineare Algebra (Modul 61112)}
\end{center}

\vspace{0.8em}

Diese Lektion löst das \textbf{Normalformproblem}: Jede komplexe quadratische Matrix ist ähnlich zu
einer Matrix in \textbf{Jordan'scher Normalform}. Der Weg führt über \textbf{nilpotente Endomorphismen}
(deren Normalformtheorie der schwierigste Schritt ist) und \textbf{Haupträume}.

Dies ist die anspruchsvollste Lektion des Moduls -- bleiben Sie am Ball. Bevor Sie sich mit
der allgemeinen Jordan-Form beschäftigen, sollten Sie zunächst die nilpotente Normalform
gründlich verstehen.

% -----------------------------------------------------------------------
\section*{Struktur der Lektion 6}

\begin{center}
\begin{tikzpicture}[node distance=0.55cm]
  \node[box] (inv) {6.1 Invariante Unterräume\\$U_\varphi(g) = \ker(g(\varphi))$,\\Zerlegung durch Minimalpolynom};
  \node[box, below=0.5cm of inv] (nil) {6.2 Nilpotente Endomorphismen\\Stufen, Filtrierung, zyklische UR,\\Rangpartition, Normalform};
  \node[box, below=0.5cm of nil] (jor) {6.3 Jordan'sche Normalform\\Jordanblöcke, Haupträume,\\Hauptsatz, Berechnung};
  \draw[arrow] (inv) -- (nil);
  \draw[arrow] (nil) -- (jor);
\end{tikzpicture}
\end{center}

\newpage

% -----------------------------------------------------------------------
\section*{Zielelement 6.1 -- Invariante Unterräume und Minimalpolynom}

\subsection*{Lerninhalte}
\begin{center}
\begin{tikzpicture}[node distance=0.55cm]
  \node[box] (ug)  {$U_\varphi(g) = \ker(g(\varphi))$\\$\varphi$-invarianter Unterraum (6.1.1)};
  \node[box, below=0.5cm of ug]  (zer) {Satz 6.1.8: Für teilerfremde $g_1, g_2$:\\$U_\varphi(g_1 g_2) = U_\varphi(g_1) \oplus U_\varphi(g_2)$};
  \node[box, below=0.5cm of zer] (dia) {Kriterium: $\varphi$ diagonal $\iff$\\$\mu_\varphi$ zerfällt in versch.\ Linearfaktoren\\(Satz 6.1.11)};
  \draw[arrow] (ug) -- (zer);
  \draw[arrow] (zer) -- (dia);
\end{tikzpicture}
\end{center}

\subsection*{Lernziele}
\begin{itemize}
  \item Den Unterraum $U_\varphi(g) = \ker(g(\varphi))$ berechnen; Invarianz nachweisen.
  \item Satz~6.1.8 anwenden: Faktorisierung des Minimalpolynoms liefert Zerlegung in invariante Unterräume.
  \item Das Kriterium~6.1.11 für Diagonalisierbarkeit über das Minimalpolynom anwenden.
\end{itemize}

\subsection*{Selbstkontrollelement 6.1}
Sei $A = \begin{pmatrix}2&0\\0&3\end{pmatrix}$. Bestimmen Sie $\mu_A$ und zerlegen Sie $\mathbb{R}^2 = U_A(X-2) \oplus U_A(X-3)$.

\newpage

% -----------------------------------------------------------------------
\section*{Zielelement 6.2 -- Nilpotente Endomorphismen}

\subsection*{Lerninhalte}
\begin{center}
\begin{tikzpicture}[node distance=0.55cm]
  \node[box] (def) {Nilpotent: $\varphi^k = 0$\\Charakterisierung (6.2.17/18)};
  \node[box, below=0.5cm of def] (fil) {Stufe eines Vektors (6.2.7)\\Filtrierung $\ker(\varphi^k)$};
  \node[box, below=0.5cm of fil] (zyk) {Zyklischer Unterraum $Z(v,\varphi)$\\Basis: $v, \varphi(v), \ldots, \varphi^{k-1}(v)$};
  \node[box, below=0.5cm of zyk] (par) {Partition, duale Partition (6.2.27--6.2.40)\\Rangpartition $p(\varphi)$};
  \node[box, below=0.5cm of par] (nf)  {Nilpotente Normalform (6.2.49/51)\\Satz: eindeutig durch Rangpartition};
  \draw[arrow] (def) -- (fil);
  \draw[arrow] (fil) -- (zyk);
  \draw[arrow] (zyk) -- (par);
  \draw[arrow] (par) -- (nf);
\end{tikzpicture}
\end{center}

\subsection*{Lernziele}
\begin{itemize}
  \item Nilpotente Endomorphismen definieren; Kriterien~6.2.17/18 kennen.
  \item Die Filtrierung $\{0\} \subsetneq \ker(\varphi) \subsetneq \ker(\varphi^2) \subsetneq \cdots$ verstehen.
  \item Zyklische Unterräume definieren und eine Basis angeben.
  \item Den Begriff der Partition und der dualen Partition erklären.
  \item Die Rangpartition $p(\varphi)$ berechnen; die nilpotente Normalform bestimmen.
  \item Zwei nilpotente Matrizen sind ähnlich $\iff$ gleiche Rangpartition (6.2.54).
\end{itemize}

\subsection*{Selbstkontrollelement 6.2}
Sei $N = \begin{pmatrix}0&1&0\\0&0&1\\0&0&0\end{pmatrix}$. Bestimmen Sie die Rangpartition und die nilpotente Normalform.

\newpage

% -----------------------------------------------------------------------
\section*{Zielelement 6.3 -- Jordan'sche Normalform}

\subsection*{Lerninhalte}
\begin{center}
\begin{tikzpicture}[node distance=0.55cm]
  \node[box] (jb)  {Jordanblock $J_k(\lambda)$ der Größe $k$\\Jordankästchen};
  \node[box, below=0.5cm of jb]  (jnf) {Jordan'sche Normalform (6.3.5):\\blockdiagonal aus Jordanblöcken};
  \node[box, below=0.5cm of jnf] (hr)  {Hauptraum $H_\varphi(\lambda) = U_\varphi((X-\lambda)^n)$\\Zerlegung $V = \bigoplus_\lambda H_\varphi(\lambda)$};
  \node[box, below=0.5cm of hr]  (hs)  {Hauptsatz (6.3.7/8):\\JNF existiert $\iff$ $\chi_\varphi$ zerfällt in LF};
  \node[box, below=0.5cm of hs]  (ber) {Berechnung der JNF\\(6.3.19): Rangpartition pro Eigenwert};
  \draw[arrow] (jb) -- (jnf);
  \draw[arrow] (jnf) -- (hr);
  \draw[arrow] (hr) -- (hs);
  \draw[arrow] (hs) -- (ber);
\end{tikzpicture}
\end{center}

\subsection*{Lernziele}
\begin{itemize}
  \item Jordanblöcke und Jordankästchen definieren; eine Matrix in JNF erkennen.
  \item Haupträume $H_\varphi(\lambda)$ berechnen; Zerlegung $V = \bigoplus H_\varphi(\lambda)$ verstehen.
  \item Den Hauptsatz~6.3.7 kennen: JNF existiert $\iff$ $\chi_\varphi$ zerfällt in Linearfaktoren (stets über $\mathbb{C}$).
  \item Eigenschaften aus der JNF ablesen: Eigenwerte, algebraische und geometrische Vielfachheiten.
  \item Die Jordan-Normalform einer gegebenen Matrix berechnen.
\end{itemize}

\subsection*{Selbstkontrollelement 6.3}
Bestimmen Sie die Jordan'sche Normalform von $A = \begin{pmatrix}3&1&0\\0&3&1\\0&0&3\end{pmatrix}$.

\end{document}
